% ---------------------------------------------------------
% PREAMBLE (same template as Companion X–XXIII)
% ---------------------------------------------------------
\documentclass[11pt,a4paper]{article}

\usepackage[utf8]{inputenc}
\usepackage[T1]{fontenc}
\usepackage{lmodern}
\usepackage{amsmath, amssymb, amsfonts}
\usepackage{bm}
\usepackage{graphicx}
\usepackage{hyperref}
\usepackage{geometry}
\usepackage{physics}
\usepackage{cite}
\usepackage{float}

\geometry{margin=1in}

% ---------------------------------------------------------
% TITLE
% ---------------------------------------------------------
\title{
\textbf{Companion XXIV}\\[4mm]
\textbf{The First Quasars in the Relaxation‑Driven Cyclic Cosmology  
Radiative Feedback, Ionization Fields, and Early AGN Environments}}

\author{
Michael Lehmann\\
Independent Researcher\\
\texttt{mi.lehmann@gmx.de}
}

\date{April 2026}

\begin{document}

\maketitle

% ---------------------------------------------------------
% ABSTRACT
% ---------------------------------------------------------
\begin{abstract}
The emergence of the first quasars marks a major transition in cosmic history, linking the 
formation of the earliest supermassive black holes (SMBHs) to the thermal, chemical, and 
ionization state of the surrounding intergalactic medium (IGM). Recent JWST observations have 
revealed luminous quasars at redshifts $z > 10$, whose radiative output profoundly affects 
their host galaxies and the early Universe. In the standard $\Lambda$CDM cosmology, the 
formation of such quasars requires extreme assumptions about black hole growth, feedback, and 
halo assembly.

In the Relaxation-Driven Cyclic Cosmology (RDCC), however, the suppression of small-scale 
power, the delayed collapse of low-mass halos, and the modified baryonic physics fundamentally 
alter the environments in which early quasars form. Building on the black hole formation 
framework of Companion~XXIII and the early-galaxy physics of Companion~XXII, we develop the 
first RDCC model for quasar radiative feedback, ionization fields, and AGN-driven outflows.

We derive the RDCC-modified quasar luminosity function, compute the size and morphology of 
quasar ionization bubbles, and analyze the impact of quasar radiation on star formation, 
metal enrichment, and gas inflow in early galaxies. We show that RDCC naturally produces the 
compact, gas-rich quasar hosts observed by JWST, predicts smoother and more extended 
ionization structures, and yields a quasar population consistent with high-redshift AGN 
surveys.

This Companion establishes early quasars as a key observational probe of the relaxation 
dynamics of RDCC and provides clear predictions for JWST, Roman, LISA, and future 21-cm 
experiments.
\end{abstract}

% ---------------------------------------------------------
% SECTION 1 — MOTIVATION AND OVERVIEW
% ---------------------------------------------------------
\section{Motivation and Overview}

The discovery of luminous quasars at redshifts $z > 10$ by JWST has opened a new window
into the earliest phases of cosmic structure formation. These quasars, powered by accreting
supermassive black holes (SMBHs) with masses $M_{\bullet} \gtrsim 10^{7}$–$10^{8}\,M_{\odot}$,
represent some of the most extreme objects in the early Universe. Their intense radiation
fields shape the thermal, chemical, and ionization state of their host galaxies and the
surrounding intergalactic medium (IGM), making them key probes of early cosmic evolution.

In the standard $\Lambda$CDM cosmology, the formation of such luminous quasars at early
times poses several challenges:
\begin{itemize}
\item rapid SMBH assembly requires sustained Eddington or super-Eddington accretion,
\item host galaxies must supply large inflow rates despite strong feedback,
\item ionization bubbles must grow rapidly in a dense IGM,
\item compact, gas-rich hosts are difficult to reconcile with early feedback.
\end{itemize}

These tensions suggest that early quasar formation may proceed differently from the
standard scenario.

% ---------------------------------------------------------
\subsection{RDCC Perspective}

The Relaxation-Driven Cyclic Cosmology (RDCC) modifies early structure formation through
the scale-dependent relaxation rate $\alpha(a)$, which suppresses small-scale power and delays
the collapse of low-mass halos. As shown in \textbf{Companion~XXII} and \textbf{Companion~XXIII},
this leads to:
\begin{itemize}
\item smoother early galaxy assembly,
\item reduced fragmentation and enhanced gas reservoirs,
\item more massive black hole seeds,
\item enhanced Eddington and super-Eddington accretion,
\item reduced radiative and mechanical feedback coupling.
\end{itemize}

These effects directly influence quasar formation and evolution:
\begin{itemize}
\item quasars ignite earlier due to accelerated SMBH growth,
\item quasar hosts remain compact and gas-rich due to delayed star formation,
\item ionization bubbles become larger and smoother due to reduced small-scale structure,
\item quasar radiation couples less efficiently to the surrounding gas, enabling sustained inflow.
\end{itemize}

Thus, RDCC provides a natural pathway to forming luminous quasars at $z > 10$ without
fine-tuned initial conditions.

% ---------------------------------------------------------
\subsection{Goals of This Companion}

The goals of this Companion are:
\begin{itemize}
\item to derive the RDCC-modified quasar luminosity function (QLF),
\item to compute the size and morphology of quasar ionization bubbles,
\item to analyze radiative and mechanical quasar feedback in RDCC halos,
\item to model the impact of quasar radiation on early galaxy evolution,
\item to compare RDCC predictions with JWST, Roman, LISA, and 21-cm observations.
\end{itemize}

These results build directly on the black hole formation framework of
\textbf{Companion~XXIII} and the early-galaxy physics of \textbf{Companion~XXII}.

% ---------------------------------------------------------
\subsection{Structure of This Companion}

The structure of this Companion is as follows:
\begin{itemize}
\item Section~2 derives the RDCC-modified quasar luminosity function.
\item Section~3 analyzes quasar radiative feedback and ionization fields.
\item Section~4 develops the RDCC quasar host galaxy model.
\item Section~5 compares RDCC predictions with JWST and high-$z$ AGN.
\item Section~6 summarizes observational signatures and provides concluding remarks.
\end{itemize}

Together, these results establish early quasars as a powerful observational probe of the
relaxation dynamics of RDCC.

% ---------------------------------------------------------
% SECTION 2 — THE RDCC QUASAR LUMINOSITY FUNCTION
% ---------------------------------------------------------
\section{The RDCC Quasar Luminosity Function}

The quasar luminosity function (QLF) encodes the abundance of accreting supermassive
black holes (SMBHs) as a function of luminosity and redshift. In the standard
$\Lambda$CDM model, the QLF at $z > 8$ is expected to decline steeply due to the scarcity
of massive halos, the difficulty of sustaining high accretion rates, and the strong radiative
feedback that limits gas inflow. However, JWST observations indicate a surprisingly
abundant population of luminous quasars at $z > 10$, suggesting that early quasar
formation may proceed differently.

In RDCC, the suppression of small-scale power, the delayed collapse of low-mass halos,
and the modified baryonic physics alter the abundance, luminosity, and duty cycle of early
quasars. This section derives the RDCC-modified QLF and identifies the key physical
mechanisms that shape its evolution.

% ---------------------------------------------------------
\subsection{From Black Hole Mass to Luminosity}

The bolometric luminosity of a quasar is
\begin{equation}
L
=
\epsilon\,
\dot{M}_{\bullet}\,
c^{2},
\end{equation}
where $\epsilon$ is the radiative efficiency.

Using the RDCC-modified accretion rate from Companion~XXIII,
\begin{equation}
\dot{M}_{\bullet,\mathrm{RDCC}}
=
f_{\mathrm{Edd,RDCC}}\,
\dot{M}_{\mathrm{Edd}}
+
f_{\mathrm{sup,RDCC}}\,
\dot{M}_{\mathrm{sup}},
\end{equation}
the luminosity becomes
\begin{equation}
L_{\mathrm{RDCC}}
=
L_{\mathrm{Edd}}
\left[
f_{\mathrm{Edd,RDCC}}
+
f_{\mathrm{sup,RDCC}}
\left(
\frac{\dot{M}_{\mathrm{sup}}}{\dot{M}_{\mathrm{Edd}}}
\right)
\right].
\end{equation}

The RDCC-modified accretion factors scale as
\begin{equation}
f_{\mathrm{Edd,RDCC}}
=
f_{\mathrm{Edd}}
\left[
1
+
\chi_{\mathrm{Edd}}
\left(
\frac{M_{\mathrm{rel}}}{M_{\bullet}}
\right)^{\beta_{\mathrm{Edd}}}
\right],
\end{equation}
\begin{equation}
f_{\mathrm{sup,RDCC}}
=
f_{\mathrm{sup}}
\left[
1
+
\chi_{\mathrm{sup}}
\left(
\frac{M_{\mathrm{rel}}}{M_{\bullet}}
\right)^{\beta_{\mathrm{sup}}}
\right],
\end{equation}
with typical values $\chi_{\mathrm{Edd}} \sim 0.1$ and $\chi_{\mathrm{sup}} \sim 0.2$.

Consequences:
\begin{itemize}
\item RDCC quasars are systematically more luminous at fixed SMBH mass,
\item super-Eddington phases are more common,
\item the bright end of the QLF is enhanced.
\end{itemize}

% ---------------------------------------------------------
\subsection{RDCC Duty Cycle}

The quasar duty cycle is the fraction of time a black hole spends in an active accretion phase.
In $\Lambda$CDM, strong feedback and limited gas supply lead to
$f_{\mathrm{duty}} \sim 0.1$–$0.2$.

In RDCC:
\begin{itemize}
\item reduced radiative feedback lowers gas expulsion,
\item smoother halo assembly increases gas inflow,
\item larger seed masses reduce the number of required accretion episodes.
\end{itemize}

Thus,
\begin{equation}
f_{\mathrm{duty,RDCC}}
=
f_{\mathrm{duty}}
\left[
1
+
\chi_{\mathrm{duty}}
\left(
\frac{M_{\mathrm{rel}}}{M_{\bullet}}
\right)^{\beta_{\mathrm{duty}}}
\right],
\end{equation}
with $\chi_{\mathrm{duty}} \sim 0.2$.

Typical RDCC values:
\begin{equation}
f_{\mathrm{duty,RDCC}}
\sim
0.3\text{--}0.5.
\end{equation}

% ---------------------------------------------------------
\subsection{Constructing the RDCC Quasar Luminosity Function}

The QLF is defined as
\begin{equation}
\Phi(L,z)
=
\frac{dn}{d\log L}.
\end{equation}

In RDCC, the QLF is determined by:
\begin{enumerate}
\item the RDCC black hole mass function (BHMF) from Companion~XXIII,
\item the RDCC-modified luminosity–mass relation,
\item the RDCC duty cycle.
\end{enumerate}

Thus,
\begin{equation}
\Phi_{\mathrm{RDCC}}(L,z)
=
f_{\mathrm{duty,RDCC}}
\int dM_{\bullet}\,
\Phi_{\bullet,\mathrm{RDCC}}(M_{\bullet},z)\,
\delta\!\left(L - L_{\mathrm{RDCC}}(M_{\bullet})\right).
\end{equation}

This expression can be evaluated analytically using the RDCC approximations of Appendix~A.

% ---------------------------------------------------------
\subsection{RDCC Predictions for the High-Redshift QLF}

RDCC predicts the following qualitative features:
\begin{itemize}
\item \textbf{Enhanced bright-end QLF:} larger seed masses and higher accretion rates increase the abundance of luminous quasars.
\item \textbf{Flatter faint-end slope:} suppressed small-scale structure reduces the number of low-luminosity AGN.
\item \textbf{Earlier QLF rise:} accelerated SMBH growth shifts the onset of quasar activity to higher redshifts.
\item \textbf{Smoother redshift evolution:} reduced fragmentation leads to more coherent quasar fueling histories.
\end{itemize}

These features match the trends observed in JWST early AGN surveys.

% ---------------------------------------------------------
\subsection{Analytical RDCC QLF Approximation}

For practical modeling, the RDCC QLF can be approximated by a modified double power law:
\begin{equation}
\Phi_{\mathrm{RDCC}}(L,z)
=
\frac{\Phi^{\ast}(z)}
{\left(L/L^{\ast}(z)\right)^{\alpha}
+
\left(L/L^{\ast}(z)\right)^{\beta}},
\end{equation}
with RDCC-modified parameters:
\begin{equation}
L^{\ast}(z)
=
L^{\ast}_{0}
\exp\!\left[
\chi_{L}
\left(
\frac{M_{\mathrm{rel}}}{M_{\bullet}}
\right)^{\beta_{L}}
\right],
\end{equation}
\begin{equation}
\Phi^{\ast}(z)
=
\Phi^{\ast}_{0}
\left[
1
+
\chi_{\Phi}
\left(
\frac{M_{\mathrm{rel}}}{M_{\bullet}}
\right)^{\beta_{\Phi}}
\right].
\end{equation}

Typical RDCC slopes:
\begin{equation}
\alpha \sim 1.2\text{--}1.4,
\qquad
\beta \sim 2.5\text{--}3.0.
\end{equation}

% ---------------------------------------------------------
\subsection{Summary}

RDCC modifies the quasar luminosity function through:
\begin{itemize}
\item enhanced luminosities at fixed black hole mass,
\item increased duty cycles,
\item accelerated SMBH growth,
\item suppressed faint-end AGN population,
\item smoother redshift evolution.
\end{itemize}

These effects set the stage for the RDCC-modified radiative feedback and ionization fields
developed in Section~3.

% ---------------------------------------------------------
% SECTION 3 — QUASAR RADIATIVE FEEDBACK AND IONIZATION FIELDS IN RDCC
% ---------------------------------------------------------
\section{Quasar Radiative Feedback and Ionization Fields in RDCC}

Quasar radiation profoundly influences the thermal, chemical, and ionization state of the
interstellar medium (ISM) and intergalactic medium (IGM). In the standard $\Lambda$CDM
model, strong radiative feedback can suppress star formation, expel gas from galaxies, and
limit the growth of supermassive black holes (SMBHs). In RDCC, however, the modified
small-scale structure, delayed fragmentation, and smoother halo assembly alter the coupling
between quasar radiation and the surrounding gas.

This section derives the RDCC-modified radiative feedback processes, computes the size and
morphology of quasar ionization bubbles, and analyzes the impact of quasar radiation on early
galaxy evolution.

% ---------------------------------------------------------
\subsection{Quasar Spectral Energy Distribution}

The quasar spectral energy distribution (SED) is modeled as
\begin{equation}
L_{\nu}
=
L_{0}
\left(
\frac{\nu}{\nu_{0}}
\right)^{-\alpha_{\nu}},
\end{equation}
with typical slopes
\begin{equation}
\alpha_{\mathrm{UV}} \sim 1.5,
\qquad
\alpha_{\mathrm{X}} \sim 0.9.
\end{equation}

In RDCC, the SED shape is unchanged, but the luminosity normalization is enhanced due to:
\begin{itemize}
\item larger seed masses,
\item higher Eddington ratios,
\item more common super-Eddington phases.
\end{itemize}

We write
\begin{equation}
L_{0,\mathrm{RDCC}}
=
L_{0}
\left[
1
+
\chi_{L}
\left(
\frac{M_{\mathrm{rel}}}{M_{\bullet}}
\right)^{\beta_{L}}
\right],
\end{equation}
with $\chi_{L} \sim 0.1$.

% ---------------------------------------------------------
\subsection{Ionizing Photon Production}

The ionizing photon rate is
\begin{equation}
\dot{N}_{\gamma}
=
\int_{\nu_{\mathrm{ion}}}^{\infty}
\frac{L_{\nu}}{h\nu}\,d\nu.
\end{equation}

Using the RDCC-modified luminosity,
\begin{equation}
\dot{N}_{\gamma,\mathrm{RDCC}}
=
\dot{N}_{\gamma}
\left[
1
+
\chi_{\gamma}
\left(
\frac{M_{\mathrm{rel}}}{M_{\bullet}}
\right)^{\beta_{\gamma}}
\right],
\end{equation}
with $\chi_{\gamma} \sim 0.1$.

Consequences:
\begin{itemize}
\item RDCC quasars produce more ionizing photons at fixed SMBH mass,
\item ionization bubbles grow faster,
\item the UV background is enhanced at early times.
\end{itemize}

% ---------------------------------------------------------
\subsection{Growth of Quasar Ionization Bubbles}

The radius of the ionized region (Strömgren sphere) is
\begin{equation}
R_{\mathrm{ion}}
=
\left[
\frac{3\,\dot{N}_{\gamma}}
{4\pi\,\alpha_{\mathrm{B}}\,n_{\mathrm{H}}^{2}}
\right]^{1/3}.
\end{equation}

In RDCC:
\begin{itemize}
\item the IGM density is lower due to delayed structure formation,
\item the ionizing photon rate is higher,
\item small-scale clumping is suppressed.
\end{itemize}

Thus,
\begin{equation}
R_{\mathrm{ion,RDCC}}
=
R_{\mathrm{ion}}
\left[
1
+
\chi_{R}
\left(
\frac{k}{k_{\mathrm{rel}}}
\right)^{\alpha_{R}}
\right],
\end{equation}
with $\chi_{R} \sim 0.2$.

Typical RDCC values:
\begin{equation}
R_{\mathrm{ion,RDCC}}
\sim
(1.5\text{--}2)\,R_{\mathrm{ion}}.
\end{equation}

% ---------------------------------------------------------
\subsection{X-ray Heating of the IGM}

X-rays from quasars heat the IGM through photoionization and secondary electrons. The
heating rate is
\begin{equation}
H_{\mathrm{X}}
=
\int_{\nu_{\mathrm{X}}}^{\infty}
\frac{L_{\nu}}{4\pi r^{2}}
\sigma_{\mathrm{HI}}(\nu)
\left(h\nu - h\nu_{\mathrm{ion}}\right)
d\nu.
\end{equation}

In RDCC:
\begin{itemize}
\item lower gas densities reduce cooling,
\item smoother density fields reduce shadowing,
\item enhanced luminosities increase heating.
\end{itemize}

Thus,
\begin{equation}
H_{\mathrm{X,RDCC}}
=
H_{\mathrm{X}}
\left[
1
+
\chi_{\mathrm{X}}
\left(
\frac{k}{k_{\mathrm{rel}}}
\right)^{\alpha_{\mathrm{X}}}
\right],
\end{equation}
with $\chi_{\mathrm{X}} \sim 0.2$.

% ---------------------------------------------------------
\subsection{Radiative Feedback on Star Formation}

Quasar radiation suppresses star formation by:
\begin{itemize}
\item heating the ISM,
\item ionizing dense gas,
\item driving outflows,
\item reducing molecular hydrogen formation.
\end{itemize}

In RDCC:
\begin{itemize}
\item delayed fragmentation increases gas reservoirs,
\item reduced small-scale clumping lowers feedback coupling,
\item smoother inflow stabilizes gas supply.
\end{itemize}

Thus, the effective star formation suppression factor becomes
\begin{equation}
SFR_{\mathrm{RDCC}}
=
SFR
\left[
1
-
\chi_{\mathrm{SF}}
\left(
\frac{M_{\mathrm{rel}}}{M_{\bullet}}
\right)^{\beta_{\mathrm{SF}}}
\right],
\end{equation}
with $\chi_{\mathrm{SF}} \sim 0.1$.

% ---------------------------------------------------------
\subsection{Mechanical Feedback: Quasar-Driven Outflows}

Quasar winds and jets inject mechanical energy into the ISM. The mechanical power is
\begin{equation}
\dot{E}_{\mathrm{mech}}
=
\epsilon_{\mathrm{mech}}\,
\dot{M}_{\bullet}\,c^{2}.
\end{equation}

In RDCC:
\begin{itemize}
\item lower gas densities reduce coupling efficiency,
\item smoother inflow stabilizes accretion disks,
\item larger seeds produce deeper potential wells.
\end{itemize}

Thus,
\begin{equation}
\dot{E}_{\mathrm{mech,RDCC}}
=
\dot{E}_{\mathrm{mech}}
\left[
1
-
\chi_{\mathrm{mech}}
\left(
\frac{M_{\mathrm{rel}}}{M_{\bullet}}
\right)^{\beta_{\mathrm{mech}}}
\right],
\end{equation}
with $\chi_{\mathrm{mech}} \sim 0.1$.

% ---------------------------------------------------------
\subsection{Summary}

RDCC modifies quasar radiative feedback through:
\begin{itemize}
\item enhanced ionizing photon production,
\item larger and smoother ionization bubbles,
\item stronger X-ray heating of the IGM,
\item reduced star formation suppression,
\item weaker mechanical feedback coupling.
\end{itemize}

These effects enable luminous quasars to form earlier and sustain high accretion rates,
setting the stage for the RDCC quasar host galaxy model developed in Section~4.

% ---------------------------------------------------------
% SECTION 4 — QUASAR HOST GALAXIES IN RDCC
% ---------------------------------------------------------
\section{Quasar Host Galaxies in RDCC}

Quasar host galaxies provide the physical environment in which supermassive black holes
(SMBHs) grow, accrete, and interact with their surroundings. Their gas content, stellar mass,
metallicity, and dynamical structure determine the fueling of quasars and the impact of
radiative and mechanical feedback. In the standard $\Lambda$CDM cosmology, early quasar
hosts are expected to be compact, rapidly star-forming, and strongly affected by feedback.
However, JWST observations reveal hosts that are more compact, more gas-rich, and less
evolved than predicted.

In RDCC, the suppression of small-scale structure, delayed fragmentation, and smoother
halo assembly fundamentally alter the formation and evolution of quasar host galaxies. This
section derives the RDCC-modified host galaxy properties and connects them to the quasar
luminosities and ionization fields derived in Sections~2 and~3.

% ---------------------------------------------------------
\subsection{Gas Content and Inflow Rates}

The gas mass of a quasar host galaxy is
\begin{equation}
M_{\mathrm{gas}}
=
f_{b}\,M_{\mathrm{halo}}
-
M_{\star},
\end{equation}
where $f_{b}$ is the cosmic baryon fraction.

In RDCC:
\begin{itemize}
\item delayed star formation reduces $M_{\star}$,
\item smoother halo assembly increases gas retention,
\item reduced fragmentation increases the cold gas fraction.
\end{itemize}

Thus,
\begin{equation}
M_{\mathrm{gas,RDCC}}
=
M_{\mathrm{gas}}
\left[
1
+
\chi_{\mathrm{gas}}
\left(
\frac{M_{\mathrm{rel}}}{M_{\mathrm{halo}}}
\right)^{\beta_{\mathrm{gas}}}
\right],
\end{equation}
with $\chi_{\mathrm{gas}} \sim 0.2$.

The gas inflow rate onto the central SMBH is
\begin{equation}
\dot{M}_{\mathrm{in}}
=
\frac{M_{\mathrm{gas}}}{t_{\mathrm{dyn}}}.
\end{equation}

In RDCC:
\begin{equation}
\dot{M}_{\mathrm{in,RDCC}}
=
\dot{M}_{\mathrm{in}}
\left[
1
+
\chi_{\mathrm{in}}
\left(
\frac{M_{\mathrm{rel}}}{M_{\mathrm{halo}}}
\right)^{\beta_{\mathrm{in}}}
\right],
\end{equation}
with $\chi_{\mathrm{in}} \sim 0.1$.

% ---------------------------------------------------------
\subsection{Stellar Mass and Star Formation}

The stellar mass of a quasar host is determined by the integrated star formation history.
In $\Lambda$CDM, early quasar hosts are expected to be rapidly star-forming, but JWST finds
hosts with:
\begin{itemize}
\item moderate stellar masses ($M_{\star} \sim 10^{8}$–$10^{9}\,M_{\odot}$),
\item suppressed star formation rates,
\item high gas fractions.
\end{itemize}

In RDCC:
\begin{itemize}
\item delayed fragmentation reduces early star formation,
\item quasar radiation couples less efficiently to the ISM,
\item smoother inflow stabilizes gas supply.
\end{itemize}

Thus, the star formation rate becomes
\begin{equation}
\mathrm{SFR}_{\mathrm{RDCC}}
=
\mathrm{SFR}
\left[
1
-
\chi_{\mathrm{SF}}
\left(
\frac{M_{\mathrm{rel}}}{M_{\mathrm{halo}}}
\right)^{\beta_{\mathrm{SF}}}
\right],
\end{equation}
with $\chi_{\mathrm{SF}} \sim 0.1$.

% ---------------------------------------------------------
\subsection{Metallicity and Chemical Enrichment}

Quasar host metallicities are shaped by:
\begin{itemize}
\item Pop~III enrichment,
\item stellar feedback,
\item inflow of pristine gas,
\item quasar-driven outflows.
\end{itemize}

In RDCC:
\begin{itemize}
\item delayed Pop~III formation delays enrichment,
\item reduced fragmentation lowers early metal production,
\item smoother inflow dilutes metallicity,
\item weaker mechanical feedback reduces metal expulsion.
\end{itemize}

Thus,
\begin{equation}
Z_{\mathrm{RDCC}}
=
Z
\left[
1
-
\chi_{Z}
\left(
\frac{M_{\mathrm{rel}}}{M_{\mathrm{halo}}}
\right)^{\beta_{Z}}
\right],
\end{equation}
with $\chi_{Z} \sim 0.1$.

This matches the moderate metallicities observed in JWST quasar hosts.

% ---------------------------------------------------------
\subsection{Morphology and Compactness}

JWST finds that early quasar hosts are:
\begin{itemize}
\item extremely compact ($R_{\mathrm{eff}} \sim 100$–$300$ pc),
\item centrally concentrated,
\item gas-dominated.
\end{itemize}

In RDCC:
\begin{itemize}
\item suppressed small-scale structure reduces fragmentation,
\item delayed star formation keeps gas centrally concentrated,
\item smoother inflow maintains compact morphologies.
\end{itemize}

Thus,
\begin{equation}
R_{\mathrm{eff,RDCC}}
=
R_{\mathrm{eff}}
\left[
1
-
\chi_{R}
\left(
\frac{M_{\mathrm{rel}}}{M_{\mathrm{halo}}}
\right)^{\beta_{R}}
\right],
\end{equation}
with $\chi_{R} \sim 0.1$.

% ---------------------------------------------------------
\subsection{Black Hole–Host Scaling Relations}

In $\Lambda$CDM, the $M_{\bullet}$–$M_{\star}$ relation is expected to hold at high redshift.
JWST finds:
\begin{itemize}
\item overmassive black holes relative to their hosts,
\item large scatter in $M_{\bullet}/M_{\star}$,
\item weak correlation at $z > 8$.
\end{itemize}

In RDCC:
\begin{itemize}
\item larger seed masses increase $M_{\bullet}$,
\item delayed star formation reduces $M_{\star}$,
\item smoother inflow stabilizes SMBH growth.
\end{itemize}

Thus,
\begin{equation}
\left(
\frac{M_{\bullet}}{M_{\star}}
\right)_{\mathrm{RDCC}}
\sim
(5\text{--}10)
\left(
\frac{M_{\bullet}}{M_{\star}}
\right)_{\Lambda\mathrm{CDM}}.
\end{equation}

This matches JWST observations.

% ---------------------------------------------------------
\subsection{Summary}

RDCC modifies quasar host galaxies through:
\begin{itemize}
\item enhanced gas fractions,
\item reduced star formation rates,
\item delayed chemical enrichment,
\item compact, centrally concentrated morphologies,
\item elevated $M_{\bullet}/M_{\star}$ ratios.
\end{itemize}

These properties match the compact, gas-rich, moderately evolved quasar hosts observed by
JWST and set the stage for the observational comparison in Section~5.

% ---------------------------------------------------------
% SECTION 5 — COMPARISON WITH JWST AND HIGH-z AGN
% ---------------------------------------------------------
\section{Comparison with JWST and High-\texorpdfstring{$z$}{z} AGN}

The emergence of luminous quasars at $z > 10$ in JWST observations provides a stringent
test for any cosmological model of early structure formation. Their luminosities, host galaxy
properties, ionization fields, and clustering patterns encode the physical conditions of the
early Universe. In this section, we compare the predictions of the Relaxation-Driven Cyclic
Cosmology (RDCC) with current observational constraints from JWST, ALMA, HST legacy
fields, and high-redshift AGN surveys.

% ---------------------------------------------------------
\subsection{Quasar Luminosities and Abundances}

JWST has identified quasars at $z > 10$ with bolometric luminosities
\begin{equation}
L_{\mathrm{bol}}
\sim
10^{45}\text{--}10^{46}\,\mathrm{erg\,s^{-1}},
\end{equation}
corresponding to black hole masses
$M_{\bullet} \sim 10^{7}$–$10^{8}\,M_{\odot}$.

In $\Lambda$CDM, producing such objects requires:
\begin{itemize}
\item sustained Eddington accretion for $\gtrsim 500$ Myr,
\item rare, massive DCBH seeds,
\item extreme duty cycles ($f_{\mathrm{duty}} \sim 1$).
\end{itemize}

RDCC predicts:
\begin{itemize}
\item larger seed masses ($10^{3}$–$10^{6}\,M_{\odot}$),
\item enhanced Eddington ratios ($f_{\mathrm{Edd,RDCC}} \sim 1.1$–$1.3$),
\item more common super-Eddington phases ($f_{\mathrm{sup,RDCC}} \sim 2$–$5$),
\item higher duty cycles ($f_{\mathrm{duty,RDCC}} \sim 0.3$–$0.5$).
\end{itemize}

\textbf{Key agreement:}  
RDCC naturally produces the observed abundance of luminous quasars at $z > 10$.

% ---------------------------------------------------------
\subsection{Quasar Host Galaxies}

JWST finds that early quasar hosts are:
\begin{itemize}
\item compact ($R_{\mathrm{eff}} \sim 100$–$300$ pc),
\item gas-rich,
\item moderately massive ($M_{\star} \sim 10^{8}$–$10^{9}\,M_{\odot}$),
\item metal-poor.
\end{itemize}

RDCC predicts:
\begin{itemize}
\item delayed star formation (Companion~XXII),
\item enhanced gas fractions (Section~4),
\item reduced metallicities due to delayed enrichment,
\item compact morphologies due to suppressed fragmentation.
\end{itemize}

\textbf{Key agreement:}  
RDCC matches the compact, gas-rich, moderately evolved quasar hosts observed by JWST.

% ---------------------------------------------------------
\subsection{Ionization Bubbles and Radiative Fields}

JWST and ALMA observations indicate:
\begin{itemize}
\item extended ionization regions around early quasars,
\item smooth ionization morphologies,
\item strong X-ray heating signatures.
\end{itemize}

RDCC predicts:
\begin{itemize}
\item larger ionization bubbles due to enhanced ionizing photon production,
\item smoother ionization fronts due to suppressed small-scale structure,
\item stronger X-ray heating due to enhanced luminosities.
\end{itemize}

\textbf{Key agreement:}  
RDCC reproduces the size and morphology of early quasar ionization bubbles.

% ---------------------------------------------------------
\subsection{Clustering and Host Halo Masses}

High-redshift AGN show:
\begin{itemize}
\item strong clustering,
\item high bias values,
\item association with massive halos.
\end{itemize}

RDCC predicts:
\begin{itemize}
\item enhanced galaxy and quasar bias (Companion~XXII),
\item suppressed low-mass halo population,
\item more massive host halos at early times.
\end{itemize}

\textbf{Key agreement:}  
RDCC naturally explains the strong clustering of early quasars.

% ---------------------------------------------------------
\subsection{Metallicity and Chemical Enrichment}

Observations indicate:
\begin{itemize}
\item moderate metallicities in quasar hosts,
\item delayed but rapid enrichment,
\item significant sightline variation.
\end{itemize}

RDCC predicts:
\begin{itemize}
\item delayed Pop~III enrichment (Companion~XXII),
\item reduced early metal production,
\item smoother inflow of pristine gas,
\item weaker mechanical feedback (Section~3).
\end{itemize}

\textbf{Key agreement:}  
RDCC matches the metallicity patterns observed in high-$z$ quasar environments.

% ---------------------------------------------------------
\subsection{LISA-Relevant Merger Rates}

LISA is sensitive to mergers with masses
\begin{equation}
M_{\bullet}
\sim
10^{3}\text{--}10^{6}\,M_{\odot}.
\end{equation}

RDCC predicts:
\begin{itemize}
\item enhanced major merger rates (Companion~XXIII),
\item more mergers with mass ratios $q > 0.1$,
\item a detectable high-$z$ merger population,
\item characteristic signatures of the relaxation scale $k_{\mathrm{rel}}$.
\end{itemize}

\textbf{Key prediction:}  
LISA can directly test the RDCC quasar assembly scenario.

% ---------------------------------------------------------
\subsection{Summary}

RDCC provides a coherent explanation for:
\begin{itemize}
\item the abundance and luminosities of $z > 10$ quasars,
\item the compact, gas-rich host galaxies observed by JWST,
\item the size and morphology of quasar ionization bubbles,
\item the strong clustering of early AGN,
\item the metallicity patterns in early quasar environments,
\item and the expected LISA merger rates.
\end{itemize}

These agreements demonstrate that early quasars are a powerful observational probe of the
relaxation dynamics of RDCC.

% ---------------------------------------------------------
% SECTION 6 — CONCLUSION
% ---------------------------------------------------------
\section{Conclusion}

In this Companion we have developed the first comprehensive model of quasar formation,
radiative feedback, and host galaxy evolution in the Relaxation-Driven Cyclic Cosmology
(RDCC). Building on the early-galaxy framework of Companion~XXII and the black hole
formation and growth model of Companion~XXIII, we have shown that the relaxation-modified
dynamics of RDCC fundamentally reshape the environments in which the first quasars ignite.

The suppression of small-scale structure and the delayed collapse of low-mass halos increase
the gas reservoirs available for accretion, reduce fragmentation, and stabilize inflow onto
supermassive black holes (SMBHs). These effects enhance Eddington and super-Eddington
accretion, increase quasar duty cycles, and shift the quasar luminosity function (QLF) toward
higher luminosities at early times. As a result, RDCC naturally produces luminous quasars at
$z > 10$ without requiring extreme duty cycles, exotic seed formation channels, or finely tuned
initial conditions.

The radiative output of RDCC quasars generates larger and smoother ionization bubbles,
stronger X-ray heating, and weaker mechanical feedback coupling than in $\Lambda$CDM.
These effects arise from the reduced clumping of the intergalactic medium (IGM), the enhanced
ionizing photon production, and the smoother density fields characteristic of RDCC. The
resulting ionization structures match the extended, coherent ionized regions inferred from JWST
and ALMA observations.

Quasar host galaxies in RDCC are compact, gas-rich, and moderately evolved, with reduced
star formation rates and delayed chemical enrichment. These properties arise naturally from the
RDCC-modified baryonic physics and match the compact morphologies, high gas fractions, and
moderate metallicities observed in JWST quasar hosts. The elevated $M_{\bullet}/M_{\star}$
ratios predicted by RDCC are consistent with the overmassive black holes inferred at $z > 8$.

Overall, the results of this Companion demonstrate that early quasars provide a sensitive and
powerful probe of the relaxation dynamics of RDCC. Their luminosities, host galaxy properties,
ionization fields, and clustering patterns encode the characteristic signatures of the relaxation
scale $k_{\mathrm{rel}}$ and offer clear observational pathways for testing the RDCC framework
with JWST, Roman, LISA, and future 21-cm experiments. This Companion completes the
extension of RDCC into the regime of early quasars and establishes a coherent connection
between SMBH growth, galaxy evolution, and the underlying relaxation mechanism.

% ---------------------------------------------------------
% APPENDIX A — ANALYTICAL RDCC REIONIZATION AND VISIBILITY APPROXIMATIONS
% ---------------------------------------------------------
\appendix
\section*{Appendix A: Analytical RDCC Reionization and Visibility Approximations}
\addcontentsline{toc}{section}{Appendix A: Analytical RDCC Reionization and Visibility Approximations}

This appendix provides compact analytical expressions for the reionization visibility function,
optical depth, and RDCC-modified large-scale potentials. These approximations complement
the numerical results of the main text and provide a transparent link between the relaxation
mechanism and the reionization signatures discussed in Sections~2–5.

% ---------------------------------------------------------
\subsection*{A.1 Optical Depth and Visibility Function}

The optical depth is defined as
\begin{equation}
\tau(z)
=
\int_{0}^{z}
dz'\,
\frac{c\,\sigma_{T}\,n_{e}(z')}{(1+z')H(z')}.
\end{equation}

The visibility function is
\begin{equation}
g(z)
=
\frac{d\tau}{dz}\,
e^{-\tau(z)}.
\end{equation}

In RDCC, the background expansion is unchanged, but the large-scale potentials entering the
ionization history are modified by the relaxation-suppressed growth rate (Companions~XVI–XVIII).

A compact RDCC approximation is
\begin{equation}
g_{\mathrm{RDCC}}(z)
=
g(z)
\left[
1
+
\chi_{g}
\left(
\frac{k_{\mathrm{rel}}}{k_{\mathrm{rei}}(z)}
\right)^{\alpha_{g}}
\right],
\end{equation}
where $k_{\mathrm{rei}}(z)$ is the characteristic scale of reionization-driven potential evolution.

This captures the enhanced large-scale coherence of the reionization visibility kernel.

% ---------------------------------------------------------
\subsection*{A.2 RDCC Modification of Large-Scale Potentials}

The large-scale Newtonian potential evolves as
\begin{equation}
\Phi(k,z)
\propto
D(z)\,T(k),
\end{equation}
where $D(z)$ is the growth factor and $T(k)$ the transfer function.

In RDCC, the relaxation mechanism modifies the growth factor as
\begin{equation}
D_{\mathrm{RDCC}}(z)
=
D(z)
\left[
1
-
\chi_{D}
\left(
\frac{k}{k_{\mathrm{rel}}}
\right)^{\alpha_{D}}
\right].
\end{equation}

Thus, the reionization-scale potential becomes
\begin{equation}
\Phi_{\mathrm{RDCC}}(k,z_{\mathrm{rei}})
=
\Phi(k,z_{\mathrm{rei}})
\left[
1
-
\chi_{\Phi}
\left(
\frac{k}{k_{\mathrm{rel}}}
\right)^{\alpha_{\Phi}}
\right].
\end{equation}

This predicts smoother large-scale potentials during reionization.

% ---------------------------------------------------------
\subsection*{A.3 Analytical Approximation for the RDCC Reionization Bump}

The reionization bump in the CMB polarization spectrum is sourced by the quadrupole at
$z \sim 6$–$10$. In RDCC, the quadrupole is modified by the relaxation-suppressed potential
evolution.

A compact approximation is
\begin{equation}
C_{\ell,\mathrm{RDCC}}^{EE,\mathrm{rei}}
\simeq
C_{\ell}^{EE,\mathrm{rei}}
\left[
1
+
\chi_{\mathrm{rei}}
\left(
\frac{\ell_{\mathrm{rel}}}{\ell}
\right)^{\alpha_{\mathrm{rei}}}
\right],
\end{equation}
with $\ell_{\mathrm{rel}} \simeq k_{\mathrm{rel}}\chi_{\mathrm{rei}}$.

This captures:
\begin{itemize}
\item enhanced large-scale polarization,
\item smoother reionization bump,
\item reduced small-scale reionization-induced fluctuations.
\end{itemize}

% ---------------------------------------------------------
\subsection*{A.4 RDCC Visibility Kernel Approximation}

The visibility kernel can be approximated as a log-normal distribution:
\begin{equation}
g(z)
\simeq
A
\exp\!\left[
-\frac{(\ln z - \ln z_{p})^{2}}{2\sigma_{z}^{2}}
\right].
\end{equation}

In RDCC:
\begin{equation}
g_{\mathrm{RDCC}}(z)
\simeq
g(z)
\left[
1
+
\chi_{v}
\left(
\frac{z}{z_{\mathrm{rel}}}
\right)^{-\alpha_{v}}
\right],
\end{equation}
with $z_{\mathrm{rel}}$ defined by $k_{\mathrm{rel}}$ through the horizon scale.

This captures the enhanced coherence of the reionization visibility window.

% ---------------------------------------------------------
\subsection*{A.5 Summary}

The analytical approximations in this appendix provide:
\begin{itemize}
\item compact expressions for RDCC-modified visibility and optical-depth kernels,
\item analytical control over the reionization bump in polarization,
\item transparent links between relaxation dynamics and reionization signatures,
\item fast semi-analytic tools for RDCC reionization modeling.
\end{itemize}

These results complement the numerical analysis of the main text and provide a practical
framework for modeling reionization in RDCC.

% ---------------------------------------------------------
% APPENDIX B — IMPLEMENTATION OF RDCC REIONIZATION IN CLASS
% ---------------------------------------------------------
\section*{Appendix B: Implementation of RDCC Reionization in CLASS}
\addcontentsline{toc}{section}{Appendix B: Implementation of RDCC Reionization in CLASS}

This appendix describes the implementation of the RDCC-modified reionization visibility
function, optical depth, and large-scale potential evolution in the CLASS Boltzmann code. The
modifications extend the RDCC growth framework of Companions~XVI–XVIII and incorporate
the analytical approximations derived in Appendix~A of this Companion.

The goal is to compute RDCC-modified reionization signatures and make them available to
CLASS and external post-processing tools.

% ---------------------------------------------------------
\subsection*{B.1 Required CLASS Modules}

RDCC reionization requires modifications in the following CLASS modules:

\begin{itemize}
\item \texttt{background.c} — access to $k_{\mathrm{rel}}(a)$ and RDCC growth functions,
\item \texttt{thermodynamics.c} — RDCC-modified visibility and optical-depth kernels,
\item \texttt{perturbations.c} — RDCC-modified large-scale potentials,
\item \texttt{input.c} — new RDCC reionization parameters,
\item \texttt{common.h} — RDCC reionization data structures,
\item \texttt{output.c} — RDCC reionization outputs.
\end{itemize}

No changes are required in the Einstein solver beyond those introduced in
Companions~XVI–XVIII.

% ---------------------------------------------------------
\subsection*{B.2 Input Parameters}

The following new input parameters must be added to \texttt{input.c}:

\begin{itemize}
\item \texttt{rdcc\_reio = yes/no} — activates RDCC reionization module,
\item \texttt{k\_rel} — relaxation scale,
\item \texttt{alpha\_rel} — suppression exponent,
\item \texttt{chi\_g} — RDCC visibility modification amplitude,
\item \texttt{chi\_Phi} — RDCC potential modification amplitude,
\item \texttt{chi\_rei} — RDCC reionization-bump amplitude.
\end{itemize}

These parameters mirror the analytical forms introduced in Appendix~A.

% ---------------------------------------------------------
\subsection*{B.3 Modifications to the Visibility Function}

The standard visibility function is computed in \texttt{thermodynamics.c} as


\[
g(z)
=
\frac{d\tau}{dz}\,e^{-\tau(z)}.
\]



In RDCC, the visibility function is modified as
\begin{equation}
g_{\mathrm{RDCC}}(z)
=
g(z)
\left[
1
+
\chi_{g}
\left(
\frac{k_{\mathrm{rel}}}{k_{\mathrm{rei}}(z)}
\right)^{\alpha_{g}}
\right],
\end{equation}
where $k_{\mathrm{rei}}(z)$ is the characteristic reionization scale.

This modification is applied after the standard CLASS visibility calculation.

% ---------------------------------------------------------
\subsection*{B.4 Modifications to the Optical Depth}

The optical depth is computed in \texttt{thermodynamics.c} as


\[
\tau(z)
=
\int dz'\,\frac{c\,\sigma_{T}\,n_{e}(z')}{(1+z')H(z')}.
\]



In RDCC, the electron density is unchanged, but the large-scale potentials entering the
ionization history are modified by the relaxation-suppressed growth rate.

A compact implementation is
\begin{equation}
\tau_{\mathrm{RDCC}}(z)
=
\tau(z)
\left[
1
+
\chi_{\tau}
\left(
\frac{k_{\mathrm{rel}}}{k_{\mathrm{rei}}(z)}
\right)^{\alpha_{\tau}}
\right].
\end{equation}

This modification is applied before computing the visibility function.

% ---------------------------------------------------------
\subsection*{B.5 Modifications to Large-Scale Potentials}

The large-scale potential is computed in \texttt{perturbations.c} as


\[
\Phi(k,z) \propto D(z)\,T(k).
\]



In RDCC:
\begin{equation}
\Phi_{\mathrm{RDCC}}(k,z)
=
\Phi(k,z)
\left[
1
-
\chi_{\Phi}
\left(
\frac{k}{k_{\mathrm{rel}}}
\right)^{\alpha_{\Phi}}
\right].
\end{equation}

This modification affects the quadrupole at reionization and therefore the reionization bump in
polarization.

% ---------------------------------------------------------
\subsection*{B.6 Reionization Bump in Polarization}

The RDCC reionization bump is implemented in \texttt{perturbations.c} as
\begin{equation}
C_{\ell,\mathrm{RDCC}}^{EE,\mathrm{rei}}
=
C_{\ell}^{EE,\mathrm{rei}}
\left[
1
+
\chi_{\mathrm{rei}}
\left(
\frac{\ell_{\mathrm{rel}}}{\ell}
\right)^{\alpha_{\mathrm{rei}}}
\right],
\end{equation}
with $\ell_{\mathrm{rel}} = k_{\mathrm{rel}}\chi_{\mathrm{rei}}$.

This modification is applied after the standard CLASS polarization source integration.

% ---------------------------------------------------------
\subsection*{B.7 Output and Post-Processing}

The RDCC reionization outputs are written to:

\begin{itemize}
\item \texttt{cl\_EE\_rei.dat},
\item \texttt{visibility\_rdcc.dat},
\item \texttt{tau\_rdcc.dat}.
\end{itemize}

These files can be used for comparison with Planck, LiteBIRD, and future polarization surveys.

% ---------------------------------------------------------
\subsection*{B.8 Summary}

The CLASS implementation of RDCC reionization requires:

\begin{itemize}
\item RDCC-modified visibility and optical-depth kernels,
\item RDCC-modified large-scale potentials,
\item reionization-bump corrections in polarization,
\item new CLASS parameters and data structures.
\end{itemize}

This implementation provides a complete numerical framework for computing RDCC
reionization signatures and enables direct comparison with current and future CMB
polarization experiments.

% ---------------------------------------------------------
% APPENDIX C — NUMERICAL SETTINGS AND VALIDATION TESTS
% ---------------------------------------------------------
\section*{Appendix C: Numerical Settings and Validation Tests}
\addcontentsline{toc}{section}{Appendix C: Numerical Settings and Validation Tests}

This appendix summarizes the numerical settings, convergence tests, and validation procedures
used to compute the RDCC reionization signatures presented in this Companion. The goal is to
ensure reproducibility and to provide guidance for implementing RDCC reionization in CLASS
and related analysis pipelines.

% ---------------------------------------------------------
\subsection*{C.1 Numerical Resolution Requirements}

Accurate computation of RDCC-modified visibility and optical-depth kernels requires sufficient
resolution in redshift and conformal time.

Recommended settings:
\begin{itemize}
\item $\Delta z \leq 0.005$ for visibility-kernel integration,
\item $\Delta \eta \leq 0.1\,\mathrm{Mpc}$ for polarization source functions,
\item $k_{\mathrm{max}} \geq 5\,h\,\mathrm{Mpc}^{-1}$ for potential evolution,
\item $\ell_{\mathrm{max}} \geq 200$ for the reionization bump in $EE$.
\end{itemize}

These settings ensure convergence of the RDCC-modified reionization spectra at the percent
level.

% ---------------------------------------------------------
\subsection*{C.2 Convergence Tests}

We perform three classes of convergence tests:

\paragraph{(1) Redshift sampling.}
The visibility function $g(z)$ is sharply peaked.  
Convergence is achieved when:


\[
\frac{\Delta C_{\ell}^{EE,\mathrm{rei}}}{C_{\ell}^{EE,\mathrm{rei}}} < 0.5\%
\quad \text{for } \ell < 20.
\]



\paragraph{(2) Potential evolution.}
The RDCC suppression factor


\[
1 - \chi_{\Phi}\left(\frac{k}{k_{\mathrm{rel}}}\right)^{\alpha_{\Phi}}
\]


requires fine sampling near $k \sim k_{\mathrm{rel}}$.  
Convergence is achieved when:


\[
\frac{\Delta \Phi_{\mathrm{RDCC}}}{\Phi_{\mathrm{RDCC}}} < 0.3\%.
\]



\paragraph{(3) Polarization source integration.}
The reionization bump is sensitive to the quadrupole at $z \sim 6$–$10$.  
Convergence is achieved when:


\[
\frac{\Delta C_{\ell}^{EE}}{C_{\ell}^{EE}} < 0.2\%
\quad \text{for } \ell < 30.
\]



% ---------------------------------------------------------
\subsection*{C.3 Validation Against $\Lambda$CDM}

RDCC reionization must reduce to $\Lambda$CDM in the limit:


\[
\chi_{g},\chi_{\Phi},\chi_{\mathrm{rei}} \rightarrow 0,
\qquad
k_{\mathrm{rel}} \rightarrow \infty.
\]



We validate:
\begin{itemize}
\item $g_{\mathrm{RDCC}}(z) \rightarrow g(z)$ at the $<0.1\%$ level,
\item $\tau_{\mathrm{RDCC}}(z) \rightarrow \tau(z)$ at the $<0.1\%$ level,
\item $C_{\ell,\mathrm{RDCC}}^{EE,\mathrm{rei}} \rightarrow C_{\ell}^{EE,\mathrm{rei}}$ at the $<0.2\%$ level,
\item the quadrupole evolution matches CLASS $\Lambda$CDM at the $<0.1\%$ level.
\end{itemize}

These checks ensure that RDCC is a controlled extension of $\Lambda$CDM.

% ---------------------------------------------------------
\subsection*{C.4 Validation of RDCC Reionization Signatures}

We validate the characteristic RDCC signatures:

\paragraph{(1) Enhanced large-scale polarization.}
For $\ell < \ell_{\mathrm{rel}}$:


\[
C_{\ell,\mathrm{RDCC}}^{EE,\mathrm{rei}}
>
C_{\ell}^{EE,\mathrm{rei}}
\quad \text{by } 10\%--25\%.
\]



\paragraph{(2) Smoother reionization bump.}
The RDCC bump is broader and less sharply peaked.

\paragraph{(3) Reduced small-scale reionization fluctuations.}
For $\ell > \ell_{\mathrm{rel}}$:


\[
C_{\ell,\mathrm{RDCC}}^{EE,\mathrm{rei}}
<
C_{\ell}^{EE,\mathrm{rei}}
\quad \text{by } 5\%--15\%.
\]



\paragraph{(4) Visibility-kernel coherence.}
The RDCC visibility kernel satisfies:


\[
g_{\mathrm{RDCC}}(z) > g(z)
\quad \text{for } z \sim 6\text{–}10.
\]



% ---------------------------------------------------------
\subsection*{C.5 Numerical Stability and Error Control}

To ensure numerical stability:

\begin{itemize}
\item enforce positivity of $g_{\mathrm{RDCC}}(z)$ and $\tau_{\mathrm{RDCC}}(z)$,
\item apply spline smoothing to RDCC suppression factors,
\item ensure smooth transitions in the quadrupole evolution,
\item use adaptive redshift sampling near the visibility peak.
\end{itemize}

These steps prevent numerical artifacts in the RDCC reionization spectra.

% ---------------------------------------------------------
\subsection*{C.6 Summary}

This appendix provides:
\begin{itemize}
\item recommended numerical settings for RDCC reionization,
\item convergence tests for redshift sampling and potential evolution,
\item validation against $\Lambda$CDM,
\item verification of RDCC reionization signatures,
\item guidelines for stable numerical implementation.
\end{itemize}

These procedures ensure that RDCC reionization predictions are accurate, reproducible, and
consistent with the analytical and numerical framework developed in this Companion.


\end{document}
